\section{Arquitectura del sistema y ruta de datos}

\subsection{Niveles de percepción multimodal}

La arquitectura de GR00T está compuesta por tres niveles de percepción visual/lingüística/de fuerza:

\begin{itemize}
    \item \textbf{Codificador visual}: procesa RGB-D y escaneos de lidar, y produce un embedding de percepción espacial
    \item \textbf{Programador de lenguaje}: analiza las instrucciones para obtener la intención de la acción, la secuencia de llamadas a herramientas y alinea el reward
    \item \textbf{Programador de acciones}: selecciona torque/force/action según las restricciones dinámicas, en un ciclo cerrado en tiempo real
\end{itemize}

\begin{knowledgebox}{Silogismo percepción-lenguaje-acción}
Estos tres niveles forman un ciclo cerrado de percepción-intención-ejecución: el módulo visual le dice al módulo de lenguaje «qué hay en la escena», el módulo de lenguaje le dice al módulo de acción «qué quiero hacer», y el módulo de acción retroalimenta el resultado y actualiza la política según el reward.
\end{knowledgebox}

\subsection{Flujo de datos y ruta de entrenamiento}

La ruta de entrenamiento típica es la siguiente:

\begin{enumerate}
    \item Recopilar datos de teleoperation/telemetry (operación humana + sensor log)
    \item Generar automáticamente el prompt de la tarea con un LLM (Eureka pipeline)
    \item Expandir la tarea en Isaac Sim y generar un synthetic replay
    \item Aplicar offline RL sobre el replay (off-policy data + behavior cloning)
    \item Aplicar sim-to-real gap alignment (domain randomization + system ID)
\end{enumerate}

\begin{importantbox}{Replay Buffer hygiene}
Mantener diversidad en el replay buffer: trajectories exitosas recientes, unas cuantas muestras fallidas, expert demonstration, LLM-generated exploration. Combinar el replay secuencial con el prioritized replay evita que el modelo sobreajuste las muestras comunes.
\end{importantbox}

\subsection{Detalles de Sim-to-Real Alignment}

La estrategia de puente de GR00T adopta una triple garantía:

\begin{enumerate}
    \item \textbf{Domain randomization}: perturba la iluminación, la textura y el coeficiente de fricción para que el modelo se adapte a escenarios físicos más amplios
    \item \textbf{System identification}: utiliza una pequeña cantidad de real-world trajectory para ajustar los parámetros de la simulación y asegurar que la dinámica corresponda
    \item \textbf{Validation loop}: ejecuta tareas simplificadas en un robot real y registra el número de casos fuera de distribución en el DevOps dashboard
\end{enumerate}

\begin{importantbox}{Sim-to-Real leash}
Cada deployment viene acompañado de un ``sim-to-real leash'': mientras el Validation loop no se apruebe, el modelo se mantiene en modo sandbox, y solo se libera cuando se detecta un tolerable drift level.
\end{importantbox}

\subsection{Resumen de la sección}

Sim-to-Real alignment requiere que domain randomization, system identification y validation loop formen un ciclo cerrado entre los tres; el leash y el guard rail garantizan una transición progresiva entre la simulación y la realidad.

